\chapter{Channel coding}\label{ch:coding}

Channel coding adds structured redundancy so that the receiver can detect and correct errors.
\Cref{fig:coding-chain} shows the coding blocks of \file{phy/coding/}. They can be combined
freely. Every encoder multiplies the \meta{code_rate} metadata by its own rate, including tail
bits and CRC overhead, so the \blk{awgn} block always refers $E_b/N_0$ to the
\emph{information} bits (\cref{eq:ebn0}).

\begin{figure}[htbp]
\centering
\begin{tikzpicture}[node distance=3mm and 3.5mm,every node/.style={font=\sffamily\scriptsize}]
  \node[blockbox,minimum width=12mm] (s) {source};
  \node[greenbox,right=of s,minimum width=12mm] (crc) {CRC};
  \node[greenbox,right=of crc,minimum width=12mm] (scr) {scrambler};
  \node[greenbox,right=of scr,minimum width=12mm] (enc) {FEC\\encoder};
  \node[greenbox,right=of enc,minimum width=12mm] (pun) {puncturer};
  \node[greenbox,right=of pun,minimum width=12mm] (il) {interleaver};
  \node[blockbox,right=of il,minimum width=12mm] (map) {mapper};
  \node[warmbox,below=7mm of map,minimum width=12mm] (ch) {channel};
  \node[blockbox,left=of ch,minimum width=12mm] (dm) {soft\\demapper};
  \node[greenbox,left=of dm,minimum width=12mm] (dil) {de-inter-\\leaver};
  \node[greenbox,left=of dil,minimum width=12mm] (dpun) {depunc-\\turer};
  \node[greenbox,left=of dpun,minimum width=12mm] (dec) {FEC\\decoder};
  \node[greenbox,left=of dec,minimum width=12mm] (dscr) {descr.};
  \node[greenbox,left=of dscr,minimum width=12mm] (dcrc) {CRC\\check};
  \draw[sig,kbits] (s)--(crc); \draw[sig,kbits] (crc)--(scr); \draw[sig,kbits] (scr)--(enc);
  \draw[sig,kbits] (enc)--(pun); \draw[sig,kbits] (pun)--(il); \draw[sig,kbits] (il)--(map);
  \draw[sig] (map)--(ch); \draw[sig] (ch)--(dm); \draw[sig,kllr] (dm)--(dil); \draw[sig,kllr] (dil)--(dpun);
  \draw[sig,kllr] (dpun)--(dec); \draw[sig,kbits] (dec)--(dscr); \draw[sig,kbits] (dscr)--(dcrc);
\end{tikzpicture}
\caption{A complete coded chain. Purple edges carry bits, pink edges LLRs. The scrambler and the
interleaver are their own inverses (\prm{inverse}\code{=true} for the interleaver).}
\label{fig:coding-chain}
\end{figure}

% ------------------------------------------------------------------
\section{Cyclic redundancy check}
A CRC with generator polynomial $g(x)$ of degree $r$ appends the remainder
$p(x)=x^r u(x)\bmod g(x)$ to the message $u(x)$; the receiver recomputes it and compares.
Because polynomial division over GF(2) is linear, the parity is a linear function of the
message, $\vect p=\vect u\mat P\bmod 2$. Row $i$ of $\mat P$ is the remainder of the monomial
$x^{k-1-i+r}$. Computing all of them once turns CRC encoding and checking of a whole batch
into a single integer matrix product.
\codeexcerpt{crc_matrix}{The CRC parity matrix, built once per frame length and cached.}
A CRC of degree $r$ detects all error bursts of length $\le r$. It misses a random error
pattern with probability $\approx2^{-r}$. \blk{crc_decoder} reports \code{bler}, the block error
rate detected by the CRC, and outputs a validity flag per frame. CRC-8, CRC-16-CCITT, LTE
CRC-24A and IEEE 802.3 CRC-32 are available.

% ------------------------------------------------------------------
\section{LFSR sequences and scrambling}
A linear feedback shift register produces $s[k]=\bigoplus_{t\in\mathcal T}s[k-t]$. When the
feedback polynomial $1+\sum_{t\in\mathcal T}x^t$ is primitive of degree $m$, the sequence is an
$m$-sequence of period $2^m-1$ with nearly ideal autocorrelation. The \blk{prbs} source
(PRBS-7 to PRBS-31, ITU-T O.150) and the additive \blk{scrambler} (default $x^7+x^4+1$, IEEE
802.11) are built on it. Since every output bit depends only on bits at least
$\min\mathcal T$ positions back, chunks of that many bits can be computed with vector operations:
\codeexcerpt{lfsr}{Vectorized LFSR generation (\file{dsp/lfsr.py}).}
Scrambling is a XOR with the sequence and is its own inverse. On LLRs, descrambling flips the
sign of the affected values. The self-synchronizing $x^{58}+x^{39}+1$ scrambler of 64B/66B is
described in \cref{sec:line-codes}.

% ------------------------------------------------------------------
\section{Interleaving}
Fading and decoders both produce errors in bursts, while most codes are designed for
independent errors. An interleaver permutes the coded bits, and the de-interleaver applies the
inverse permutation, so a burst on the channel is spread over many codewords. The
\prm{block} interleaver writes row by row into an $R\times(n/R)$ matrix and reads column by
column. The \prm{random} interleaver applies a seeded pseudo-random permutation.

% ------------------------------------------------------------------
\section{Block codes: repetition and Hamming}
The repetition code sends every bit $n$ times. Its soft decoder simply adds the $n$ LLRs. The
Hamming code $(2^m-1,2^m-1-m)$ has as parity-check matrix $\mat H$ all non-zero $m$-bit columns,
arranged so that the code is systematic. A single error in position $j$ produces the syndrome
$\vect s=\mat H\vect e\T=\vect h_j$, which identifies the position; a lookup table maps each
syndrome to the bit to flip. Decoding errors occur with $\ge2$ channel errors; the approximation
\begin{equation}
  P_b\approx\frac1n\sum_{j=2}^{n}(j+1)\binom nj p^j(1-p)^{n-j},\qquad p=\Q\Big(\sqrt{2R\,E_b/N_0}\Big),
\end{equation}
counts the $j$ channel errors plus at most one wrong correction (\code{ber_hamming_hard}).

% ------------------------------------------------------------------
\section{Convolutional codes and the Viterbi algorithm}\label{sec:conv}
\subsection{Encoder and trellis}
A rate-$1/n$ feed-forward convolutional encoder with constraint length $K$ computes $n$ output
bits per input bit,
\begin{equation}
  c_j[t]=\bigoplus_{i=0}^{K-1}g_{j,i}\,u[t-i],\qquad j=1,\dots,n,
\end{equation}
with generators given in octal: $(171,133)_8$ is the industry-standard $K=7$ code of IEEE
802.11, DVB and many others, with free distance $d_{\mathrm{free}}=10$. The encoder state is
the last $K-1$ input bits, giving $2^{K-1}$ states. Appending $K-1$ zeros (\prm{terminate})
returns the trellis to state 0.

\begin{figure}[htbp]
\centering
\begin{tikzpicture}[x=2.1cm,y=0.95cm]
  \foreach \t in {0,...,4} {
    \foreach \s/\lab in {0/00,1/01,2/10,3/11} \node[dot] (n\t\s) at (\t,-\s) {};
    \node[lbl,above] at (\t,0.25) {$t=\t$};}
  \foreach \s/\lab in {0/00,1/01,2/10,3/11} \node[lbl,left] at (-0.1,-\s) {\lab};
  \foreach \t [evaluate=\t as \u using int(\t+1)] in {0,...,3} {
    % input 0 (solid), input 1 (dashed); state = (u[t-1], u[t-2]) MSB = most recent
    \draw[accent,thin] (n\t0)--(n\u0); \draw[accent2,thin,dashed] (n\t0)--(n\u2);
    \draw[accent,thin] (n\t1)--(n\u0); \draw[accent2,thin,dashed] (n\t1)--(n\u2);
    \draw[accent,thin] (n\t2)--(n\u1); \draw[accent2,thin,dashed] (n\t2)--(n\u3);
    \draw[accent,thin] (n\t3)--(n\u1); \draw[accent2,thin,dashed] (n\t3)--(n\u3);}
  \draw[accent3,line width=1.6pt] (n00)--(n12)--(n23)--(n31)--(n40);
  \node[lbl,right] at (4.1,-0.0) {survivor};
  \draw[accent,thin] (4.6,-2.2)--(4.9,-2.2) node[right,lbl]{$u=0$};
  \draw[accent2,thin,dashed] (4.6,-2.8)--(4.9,-2.8) node[right,lbl]{$u=1$};
\end{tikzpicture}
\caption{Trellis of the $K=3$, $(7,5)_8$ code (4 states, MSB = most recent input). The Viterbi
algorithm keeps one survivor path per state and time; the green path is the decoded sequence
\code{1 1 0 0}.}
\label{fig:trellis}
\end{figure}

\subsection{Maximum-likelihood sequence decoding}
Given channel LLRs $L_{t,j}$ for the coded bits, the log-likelihood of a candidate codeword is,
up to a constant,
\begin{equation}
  \ln p(\vect y\,|\,\vect c)=\sum_{t,j}\frac{L_{t,j}}{2}(1-2c_{t,j})+\text{const},
\end{equation}
a \emph{correlation} metric. The Viterbi algorithm maximizes it over all paths of the trellis
by dynamic programming. At each time, for every state it keeps the better of the two incoming
paths (add--compare--select) and stores the decision. A traceback from state 0 (terminated
trellis) recovers the input bits. Path metrics are renormalized every step to avoid overflow.
\codeexcerpt{viterbi}{Viterbi decoding vectorized over states and frames (\code{Trellis.decode}).}
With hard decisions (bit inputs) the metric degenerates to the Hamming distance, losing about
2\,dB (\file{tests/test_physics.py::test_conv_coded_soft_vs_hard}). The asymptotic coding gain
of soft decoding is $10\log_{10}(R\,d_{\mathrm{free}})=7$\,dB for the $K=7$ code.

\subsection{Puncturing}
Deleting coded bits periodically raises the rate without a new decoder. The default pattern
\code{[1,1,1,0,0,1]} turns the rate-1/2 code into rate 3/4, as in IEEE 802.11. The
depuncturer reinserts the deleted positions as erasures ($L=0$, no information) before
decoding. When the original length is not available in the metadata, it is inferred as the
shortest mother length whose pattern keeps exactly the received number of bits.
\codeexcerpt{depuncture}{Erasure insertion (\blk{depuncturer}).}

% ------------------------------------------------------------------
\section{LDPC codes}\label{sec:ldpc}
\subsection{Structure}
A low-density parity-check code~\cite{gallager1962} is the null space of a sparse matrix $\mat H\in\{0,1\}^{m\times n}$:
$\mat H\vect c\T=\vect0$. Its Tanner graph has a variable node per code bit and a check node per
row (\cref{fig:tanner}). LinkStudio builds quasi-cyclic codes in the style of IEEE 802.11n and 5G
NR: $\mat H$ is an array of $z\times z$ circulant permutation matrices (or zero blocks), described
by a base matrix of shifts (\cref{fig:ldpc-base}).

\begin{figure}[htbp]
\centering
\begin{tikzpicture}[x=1.15cm,y=1cm]
  \foreach \i in {1,...,8} \node[circle,draw=accent,fill=soft,minimum size=6mm,inner sep=0pt,font=\scriptsize] (v\i) at (\i,0) {$c_\i$};
  \foreach \j in {1,...,4} \node[rectangle,draw=accent2,fill=softwarm,minimum size=6mm,font=\scriptsize] (c\j) at (2*\j-0.5,2) {$\oplus$};
  \foreach \j/\list in {1/{1,2,4,5},2/{2,3,6,7},3/{1,3,5,8},4/{4,6,7,8}}
    \foreach \i in \list \draw[black!50] (c\j.south)--(v\i.north);
  \draw[->,accent3,thick] (v5.north east) .. controls +(0.4,0.7) .. node[right,lbl]{$L_{v\to c}$} ($(c3.south)+(0.15,0)$);
  \node[lbl,right] at (8.5,0) {variable nodes};
  \node[lbl,right] at (8.5,2) {check nodes};
\end{tikzpicture}
\caption{A Tanner graph. Belief propagation exchanges messages along the edges: variables tell
checks what they believe, checks reply with what the other bits imply.}
\label{fig:tanner}
\end{figure}

\begin{figure}[htbp]
\centering
\begin{tikzpicture}[infoblk/.style={draw=accent,fill=soft,minimum size=5.6mm,inner sep=0pt},
  parblk/.style={draw=accent3,fill=softgreen,minimum size=5.6mm,inner sep=0pt},
  zeroblk/.style={draw=black!12,fill=white,minimum size=5.6mm,inner sep=0pt}]
  % Base matrix of the default QC-LDPC code (k_blocks=8, m_blocks=8, z=48), from QCLDPCCode.H
\node[zeroblk] at (0.00,0.00) {};
\node[zeroblk] at (0.62,0.00) {};
\node[infoblk] at (1.24,0.00) {\tiny 36};
\node[zeroblk] at (1.86,0.00) {};
\node[infoblk] at (2.48,0.00) {\tiny 12};
\node[infoblk] at (3.10,0.00) {\tiny 47};
\node[infoblk] at (3.72,0.00) {\tiny 13};
\node[zeroblk] at (4.34,0.00) {};
\node[parblk] at (4.96,0.00) {\tiny 0};
\node[zeroblk] at (5.58,0.00) {};
\node[zeroblk] at (6.20,0.00) {};
\node[zeroblk] at (6.82,0.00) {};
\node[zeroblk] at (7.44,0.00) {};
\node[zeroblk] at (8.06,0.00) {};
\node[zeroblk] at (8.68,0.00) {};
\node[zeroblk] at (9.30,0.00) {};
\node[zeroblk] at (0.00,-0.62) {};
\node[infoblk] at (0.62,-0.62) {\tiny 12};
\node[zeroblk] at (1.24,-0.62) {};
\node[zeroblk] at (1.86,-0.62) {};
\node[zeroblk] at (2.48,-0.62) {};
\node[zeroblk] at (3.10,-0.62) {};
\node[zeroblk] at (3.72,-0.62) {};
\node[infoblk] at (4.34,-0.62) {\tiny 24};
\node[parblk] at (4.96,-0.62) {\tiny 0};
\node[parblk] at (5.58,-0.62) {\tiny 0};
\node[zeroblk] at (6.20,-0.62) {};
\node[zeroblk] at (6.82,-0.62) {};
\node[zeroblk] at (7.44,-0.62) {};
\node[zeroblk] at (8.06,-0.62) {};
\node[zeroblk] at (8.68,-0.62) {};
\node[zeroblk] at (9.30,-0.62) {};
\node[infoblk] at (0.00,-1.24) {\tiny 39};
\node[infoblk] at (0.62,-1.24) {\tiny 19};
\node[zeroblk] at (1.24,-1.24) {};
\node[infoblk] at (1.86,-1.24) {\tiny 14};
\node[zeroblk] at (2.48,-1.24) {};
\node[zeroblk] at (3.10,-1.24) {};
\node[zeroblk] at (3.72,-1.24) {};
\node[infoblk] at (4.34,-1.24) {\tiny 5};
\node[zeroblk] at (4.96,-1.24) {};
\node[parblk] at (5.58,-1.24) {\tiny 0};
\node[parblk] at (6.20,-1.24) {\tiny 0};
\node[zeroblk] at (6.82,-1.24) {};
\node[zeroblk] at (7.44,-1.24) {};
\node[zeroblk] at (8.06,-1.24) {};
\node[zeroblk] at (8.68,-1.24) {};
\node[zeroblk] at (9.30,-1.24) {};
\node[infoblk] at (0.00,-1.86) {\tiny 6};
\node[zeroblk] at (0.62,-1.86) {};
\node[infoblk] at (1.24,-1.86) {\tiny 25};
\node[infoblk] at (1.86,-1.86) {\tiny 21};
\node[infoblk] at (2.48,-1.86) {\tiny 24};
\node[zeroblk] at (3.10,-1.86) {};
\node[zeroblk] at (3.72,-1.86) {};
\node[zeroblk] at (4.34,-1.86) {};
\node[zeroblk] at (4.96,-1.86) {};
\node[zeroblk] at (5.58,-1.86) {};
\node[parblk] at (6.20,-1.86) {\tiny 0};
\node[parblk] at (6.82,-1.86) {\tiny 0};
\node[zeroblk] at (7.44,-1.86) {};
\node[zeroblk] at (8.06,-1.86) {};
\node[zeroblk] at (8.68,-1.86) {};
\node[zeroblk] at (9.30,-1.86) {};
\node[zeroblk] at (0.00,-2.48) {};
\node[zeroblk] at (0.62,-2.48) {};
\node[zeroblk] at (1.24,-2.48) {};
\node[infoblk] at (1.86,-2.48) {\tiny 5};
\node[zeroblk] at (2.48,-2.48) {};
\node[zeroblk] at (3.10,-2.48) {};
\node[infoblk] at (3.72,-2.48) {\tiny 44};
\node[infoblk] at (4.34,-2.48) {\tiny 14};
\node[zeroblk] at (4.96,-2.48) {};
\node[zeroblk] at (5.58,-2.48) {};
\node[zeroblk] at (6.20,-2.48) {};
\node[parblk] at (6.82,-2.48) {\tiny 0};
\node[parblk] at (7.44,-2.48) {\tiny 0};
\node[zeroblk] at (8.06,-2.48) {};
\node[zeroblk] at (8.68,-2.48) {};
\node[zeroblk] at (9.30,-2.48) {};
\node[zeroblk] at (0.00,-3.10) {};
\node[zeroblk] at (0.62,-3.10) {};
\node[zeroblk] at (1.24,-3.10) {};
\node[zeroblk] at (1.86,-3.10) {};
\node[infoblk] at (2.48,-3.10) {\tiny 9};
\node[infoblk] at (3.10,-3.10) {\tiny 23};
\node[zeroblk] at (3.72,-3.10) {};
\node[zeroblk] at (4.34,-3.10) {};
\node[zeroblk] at (4.96,-3.10) {};
\node[zeroblk] at (5.58,-3.10) {};
\node[zeroblk] at (6.20,-3.10) {};
\node[zeroblk] at (6.82,-3.10) {};
\node[parblk] at (7.44,-3.10) {\tiny 0};
\node[parblk] at (8.06,-3.10) {\tiny 0};
\node[zeroblk] at (8.68,-3.10) {};
\node[zeroblk] at (9.30,-3.10) {};
\node[infoblk] at (0.00,-3.72) {\tiny 45};
\node[infoblk] at (0.62,-3.72) {\tiny 39};
\node[infoblk] at (1.24,-3.72) {\tiny 40};
\node[zeroblk] at (1.86,-3.72) {};
\node[zeroblk] at (2.48,-3.72) {};
\node[zeroblk] at (3.10,-3.72) {};
\node[infoblk] at (3.72,-3.72) {\tiny 25};
\node[zeroblk] at (4.34,-3.72) {};
\node[zeroblk] at (4.96,-3.72) {};
\node[zeroblk] at (5.58,-3.72) {};
\node[zeroblk] at (6.20,-3.72) {};
\node[zeroblk] at (6.82,-3.72) {};
\node[zeroblk] at (7.44,-3.72) {};
\node[parblk] at (8.06,-3.72) {\tiny 0};
\node[parblk] at (8.68,-3.72) {\tiny 0};
\node[zeroblk] at (9.30,-3.72) {};
\node[zeroblk] at (0.00,-4.34) {};
\node[zeroblk] at (0.62,-4.34) {};
\node[zeroblk] at (1.24,-4.34) {};
\node[zeroblk] at (1.86,-4.34) {};
\node[zeroblk] at (2.48,-4.34) {};
\node[infoblk] at (3.10,-4.34) {\tiny 5};
\node[zeroblk] at (3.72,-4.34) {};
\node[zeroblk] at (4.34,-4.34) {};
\node[zeroblk] at (4.96,-4.34) {};
\node[zeroblk] at (5.58,-4.34) {};
\node[zeroblk] at (6.20,-4.34) {};
\node[zeroblk] at (6.82,-4.34) {};
\node[zeroblk] at (7.44,-4.34) {};
\node[zeroblk] at (8.06,-4.34) {};
\node[parblk] at (8.68,-4.34) {\tiny 0};
\node[parblk] at (9.30,-4.34) {\tiny 0};

  \draw[decorate,decoration={brace,amplitude=5pt}] (-0.3,0.45) -- node[above=5pt,lbl]{information part $\mat H_i$ ($k_b=8$ block columns)} (4.64,0.45);
  \draw[decorate,decoration={brace,amplitude=5pt}] (4.66,0.45) -- node[above=5pt,lbl]{parity part $\mat H_p$ (dual diagonal)} (9.6,0.45);
\end{tikzpicture}
\caption{Base matrix of the default code ($k_b=m_b=8$, $z=48$, rate 1/2, $n=768$), read back from
\code{QCLDPCCode.H}. Each number is the cyclic shift of a $48\times48$ identity; empty cells are
zero blocks. Every information column has weight 3.}
\label{fig:ldpc-base}
\end{figure}

\subsection{Encoding by accumulation}
Split $\mat H=[\mat H_i\,|\,\mat H_p]$ and the codeword $\vect c=[\vect u\,|\,\vect p]$. The
parity part is block dual-diagonal: identities on the diagonal and on the sub-diagonal. The
parity equations $\mat H_i\vect u\T\oplus\mat H_p\vect p\T=\vect0$ then read, block by block,
\begin{equation}
  \vect p_0=\vect s_0,\qquad \vect p_j=\vect p_{j-1}\oplus\vect s_j,\qquad \vect s=\mat H_i\vect u\T ,
\end{equation}
so the parity is a cumulative XOR of the syndrome blocks. Encoding costs one sparse product and
needs no Gaussian elimination. The dual-diagonal structure also guarantees that $\mat H$ has
full rank.
\codeexcerpt{ldpc_encode}{Systematic QC-LDPC encoding.}

\subsection{Belief-propagation decoding}
With channel LLRs $L_v$, the flooding schedule alternates two updates.
\begin{align}
  \text{variable to check:}\quad & L_{v\to c}=L_v+\sum_{c'\in\mathcal N(v)\setminus c}L_{c'\to v},\\
  \text{check to variable:}\quad & L_{c\to v}=2\tanh^{-1}\!\prod_{v'\in\mathcal N(c)\setminus v}\tanh\frac{L_{v'\to c}}{2}
  \quad\text{(sum-product)},\\
  & L_{c\to v}\approx\alpha\prod_{v'}\sgn(L_{v'\to c})\cdot\min_{v'\in\mathcal N(c)\setminus v}|L_{v'\to c}|
  \quad\text{(normalized min-sum)}.
\end{align}
The min-sum rule only needs, for each check, the two smallest input magnitudes: the output to
the edge holding the minimum is the second minimum, and to all others the minimum. The factor
$\alpha\approx0.8$ compensates its overestimation of reliabilities. After every iteration the
hard decisions are tested against the syndrome. Frames that satisfy it stop early, and the
average number of iterations is reported as a metric.
\codeexcerpt{ldpc_minsum}{The normalized min-sum check-node update, vectorized over all checks of all frames.}
The message passing is vectorized with a padded table of the edges of every check node
(irregular row weights are padded with neutral entries) and sparse incidence matrices for the
variable-node sums.

% ------------------------------------------------------------------
\section{Polar codes}\label{sec:polar}
\subsection{Channel polarization}
Arıkan's transform~\cite{arikan2009} $\vect x=\vect u\mat G_N$, with
$\mat G_N=\begin{bsmallmatrix}1&0\\1&1\end{bsmallmatrix}^{\otimes n}$, $N=2^n$, combines $N$
uses of a channel $W$ into $N$ synthetic channels $W_N^{(i)}$ seen by the bits $u_i$ under
successive decoding. As $N$ grows they polarize: a fraction tends to noiseless channels, and the
rest to useless ones. One step splits a channel with Bhattacharyya parameter $Z$ into a worse
and a better one,
\begin{equation}
  Z(W^-)\le2Z-Z^2,\qquad Z(W^+)=Z^2 .
\end{equation}
Starting from $Z=e^{-E_s/N_0}$ (BPSK over AWGN at the \prm{design_snr_db}) and applying the
recursion $n$ times ranks the synthetic channels. The $K$ best carry information and the others
are \emph{frozen} to 0.
\codeexcerpt{polar_rel}{Reliability ranking by the Bhattacharyya recursion.}

\begin{figure}[htbp]
\centering
\begin{tikzpicture}
\begin{semilogyaxis}[width=11cm,height=5cm,xlabel={synthetic channel index $i$},ylabel={$Z(W_N^{(i)})$},
  xmin=-1,xmax=64,ymin=1e-12,ymax=2,legend pos=south west]
  \addplot[sim={accent2,mark=x,mark size=2pt}] table[x=i,y=z]{\datafile{polar_frozen}};
  \addplot[sim=accent3] table[x=i,y=z]{\datafile{polar_info}};
  \legend{frozen,information}
\end{semilogyaxis}
\end{tikzpicture}
\caption{Polarization for $N=64$, $K=32$ at a design SNR of 0\,dB (\code{reliability},
\code{frozen_mask}). Reliable channels concentrate at high indices, the frozen set at low
indices.}
\label{fig:polar}
\end{figure}

\subsection{Encoding and successive-cancellation decoding}
Without bit reversal, $\vect u=[\vect u_a,\vect u_b]$ encodes recursively as
$\vect x=[\mathrm{enc}(\vect u_a)\oplus\mathrm{enc}(\vect u_b),\ \mathrm{enc}(\vect u_b)]$,
computed in place with $\log_2N$ butterfly stages:
\codeexcerpt{polar_encode}{The polar transform in $\log_2 N$ vectorized stages.}
The SC decoder follows the same recursion on LLRs. For the halves $\vect L_1,\vect L_2$,
\begin{equation}
  f(a,b)=\sgn a\,\sgn b\,\min(|a|,|b|)\ \ (\text{left child}),\qquad
  g(a,b,\hat x)=b+(1-2\hat x)\,a\ \ (\text{right child}).
\end{equation}
The left half is decoded first; its re-encoded decisions $\hat{\vect x}_a$ enter the right half.
Subtrees made only of frozen bits return zeros immediately, and subtrees made only of
information bits (rate-1 nodes) are decided by a direct hard decision. Both shortcuts are exact.
\codeexcerpt{polar_sc}{Recursive SC decoding with rate-0 and rate-1 node shortcuts.}

% ------------------------------------------------------------------
\section{Comparison}
\begin{figure}[htbp]
\centering
\begin{tikzpicture}
\begin{semilogyaxis}[berplot,width=9.6cm,height=7.5cm,xmin=0,xmax=10,ymin=1e-6,ymax=0.5,
  legend style={at={(1.03,1)},anchor=north west}]
  \addplot[th=black!60] table[x=ebn0,y=uncoded]{\datafile{coding_theory}};
  \addplot[th=accent4,dashed] table[x=ebn0,y=hamming]{\datafile{coding_theory}};
  \addplot[sim=accent4] table[x=ebn0,y=hamming]{\datafile{coding_sim}};
  \addplot[simline={accent}] table[x=ebn0,y=conv]{\datafile{coding_sim}};
  \addplot[simline={accent3,mark=square*}] table[x=ebn0,y=polar]{\datafile{coding_sim}};
  \addplot[simline={accent2,mark=triangle*}] table[x=ebn0,y=ldpc]{\datafile{coding_sim}};
  \legend{uncoded (theory),{Hamming(7,4) hard (approx.)},{Hamming(7,4) simulated},
    conv. $K{=}7$ soft Viterbi,polar $N{=}256$ SC,QC-LDPC $n{=}768$ min-sum}
\end{semilogyaxis}
\end{tikzpicture}
\caption{QPSK over AWGN with the codes of this chapter (rate $\approx1/2$ except Hamming 4/7;
max-log LLRs; simulated with the real blocks, up to 200 errors or 20 batches per point). At
$10^{-4}$ the Hamming code gains well under 1\,dB with hard decisions, while the convolutional and
LDPC codes gain about 5 and 5.5\,dB.}
\label{fig:coding}
\end{figure}

\begin{keypoint}
Soft decisions are worth about 2\,dB, and iterative decoding of long sparse codes approaches
capacity. The rate-1/2 limit for binary-input AWGN is 0.19\,dB. The 768-bit LDPC code is about
3\,dB away at $10^{-5}$; longer codes get closer.
\end{keypoint}
